% TODO make hw stylesheet
\documentclass[10pt]{article}
\usepackage[letterpaper,top=1in,left=1in,right=1in]{geometry}
\usepackage[dvipsnames,svgnames]{xcolor}
\usepackage[shortlabels]{enumitem}
\usepackage[framemethod=TikZ]{mdframed}
\usepackage{amsmath,amssymb,amsthm}
\usepackage[colorlinks]{hyperref}
\usepackage{multicol}
\usepackage{tabularx}
\usepackage{bbm}

\hypersetup{
    colorlinks=true,
    linkcolor=blue,
    filecolor=magenta,
    urlcolor=magenta,
    pdftitle={MAT 549 HW}
}

\usepackage{mathtools}
\usepackage{thmtools}
\usepackage[textsize=footnotesize]{todonotes}
\usepackage{titling}
\usepackage{cancel}

\reversemarginpar
\mdfdefinestyle{mdbluebox}{roundcorner=10pt,innerbottommargin=9pt,
linecolor=blue,backgroundcolor=TealBlue!5,}
\declaretheoremstyle[headfont=\sffamily\bfseries\color{MidnightBlue},
mdframed={style=mdbluebox},]{thmbluebox}

\mdfdefinestyle{mdbloobox}{frametitlefont=\bfseries,innerbottommargin=8pt,
nobreak=true,backgroundcolor=TealBlue!5,linecolor=blue,}
\declaretheoremstyle[headfont=\bfseries\color{MidnightBlue},
mdframed={style=mdbloobox},headpunct={\\[3pt]},postheadspace=0pt,]{thmbloobox}

\mdfdefinestyle{mdredbox}{frametitlefont=\bfseries,innerbottommargin=8pt,
nobreak=true,backgroundcolor=Salmon!5,linecolor=RawSienna,}
\declaretheoremstyle[headfont=\bfseries\color{RawSienna},
mdframed={style=mdredbox},headpunct={\\[3pt]},postheadspace=0pt,]{thmredbox}

\mdfdefinestyle{mdgreenbox}{linecolor=ForestGreen,backgroundcolor=ForestGreen!5,
linewidth=2pt,rightline=false,leftline=true,topline=false,bottomline=false,}
\declaretheoremstyle[headfont=\bfseries\sffamily\color{ForestGreen!70!black},
mdframed={style=mdgreenbox},headpunct={ --- },]{thmgreenbox}

\mdfdefinestyle{mdorangebox}{linecolor=RedOrange,backgroundcolor=RedOrange!5,
linewidth=2pt,rightline=false,leftline=true,topline=false,bottomline=false,}
\declaretheoremstyle[headfont=\bfseries\sffamily\color{RedOrange!70!black},
mdframed={style=mdorangebox},headpunct={ --- },]{thmorangebox}

\mdfdefinestyle{mdblackbox}{linecolor=black,backgroundcolor=RedViolet!5!gray!5,
linewidth=3pt,rightline=false,leftline=true,topline=false,bottomline=false,}
\declaretheoremstyle[mdframed={style=mdblackbox}]{thmblackbox}

\declaretheoremstyle[spaceabove=6pt,spacebelow=6pt]{boldhead}
\declaretheoremstyle[spaceabove=6pt,spacebelow=6pt,headfont=\normalfont\itshape]{italichead}

\declaretheorem[style=boldhead,name=Problem]{problem}
\declaretheorem[style=boldhead,name=Problem,numbered=no]{problem*}
\declaretheorem[style=boldhead,name=Setup]{setup} % for config geo
\declaretheorem[style=boldhead,name=Example,sibling=problem]{example}
\declaretheorem[style=boldhead,name=$\bigstar$ Problem,sibling=problem]{niceprob}
\declaretheorem[style=boldhead,name=Theorem,sibling=problem]{theorem}
\declaretheorem[style=boldhead,name=Corollary,sibling=theorem]{corollary}
\declaretheorem[style=boldhead,name=Proposition,sibling=theorem]{prop}
\declaretheorem[style=boldhead,name=Lemma,numberwithin=problem]{lemma}
\declaretheorem[style=boldhead,name=Theorem,numbered=no]{theorem*}
\declaretheorem[style=boldhead,name=Lemma,numbered=no]{lemma*}
\declaretheorem[style=boldhead,name=Claim,numberwithin=problem]{claim}
\declaretheorem[style=boldhead,name=Claim,numbered=no]{claim*}
\declaretheorem[style=boldhead,name=Remark,numberwithin=problem]{remark}
\declaretheorem[style=boldhead,name=Remark,numbered=no]{remark*}
\declaretheorem[style=boldhead,name=Definition,sibling=theorem]{definition}
\declaretheorem[style=boldhead,name=Definition,numbered=no]{definition*}

\providecommand{\ol}{\overline}
\providecommand{\ceil}[1]{\left\lceil #1 \right\rceil}
\providecommand{\floor}[1]{\left\lfloor #1 \right\rfloor}
\newcommand{\dd}{\mathrm{d}}
\providecommand{\ul}{\underline}
\providecommand{\eps}{\varepsilon}
\providecommand{\half}{\frac{1}{2}}
\providecommand{\CC}{\mathbb C}
\providecommand{\FF}{\mathbb F}
\providecommand{\II}{\mathbb I}
\providecommand{\NN}{\mathbb N}
\providecommand{\QQ}{\mathbb Q}
\providecommand{\RR}{\mathbb R}
\providecommand{\ZZ}{\mathbb Z}
\providecommand{\ind}{\mathbbm 1}
\providecommand{\dg}{^\circ}
\providecommand{\ii}{\item}
\providecommand{\setdiff}{\smallsetminus}
\providecommand{\alert}[1]{{\sffamily\textbf{\textcolor{blue}{#1}}}}
\providecommand{\inv}{^{-1}}
\providecommand{\mb}[1]{\mathbf{#1}}
\newcommand{\what}{\widehat}

\newenvironment{soln}{\begin{proof}[Solution]}{\end{proof}}

\providecommand{\pow}{\operatorname{Pow}}
\providecommand{\vocab}[1]{{\textbf{\textcolor{ForestGreen}{#1}}}}
\providecommand{\lcm}{\operatorname{lcm}}
\providecommand{\abs}[1]{\left\lvert #1\right\rvert}
\providecommand{\smabs}[1]{\lvert #1\rvert}
\providecommand{\innerprod}[1]{\langle\!\langle #1\rangle\!\rangle}
\providecommand{\norm}[1]{\left\| #1\right\|}
\providecommand{\smnorm}[1]{\| #1\|}
\providecommand{\gr}{\operatorname{gr}}

\usepackage{mathrsfs}
\providecommand{\tanvec}[1]{\frac{\partial}{\partial #1}}

% place title at top right
\setlength{\droptitle}{-8ex}
\pretitle{\begin{flushright}\Large\bfseries}
\posttitle{\par\end{flushright}}
\preauthor{\begin{flushright}}
\postauthor{\end{flushright}}
\predate{\begin{flushright}}
\postdate{\end{flushright}}

\title{MAT 549 HW05}
\author{\large Aditya Pahuja}
\date{\large Due 04/14}
\begin{document}
\maketitle

\begin{problem}
    [15-1]
    Suppose $M$ is a smooth manifold that is the union of
    two orientable open submanifolds with connected intersection.
    Show that $M$ is orientable.
    Use this to give another proof that $S^n$ is orientable.
\end{problem}

\begin{soln}
    Let $U$ and $V$ be the two open submanifolds
    and let $\omega$ and $\eta$ respectively
    be orientation forms on $U$ and $V$.
    Since $U\cap V$ is connected, we can assume that
    the restrictions of $\omega$ and $\eta$ to $U\cap V$
    induce the same orientation on $U\cap V$
    by negating $\eta$ if necessary.
    Let $\{\varphi_U,\varphi_V\}$ be a partition of unity
    subordinate to the open cover $\{U,V\}$ of $M$
    (so $\varphi_U$ is zero outside $U$
    and $\varphi_V$ is zero outside $V$); then I claim that
    $\varphi_U\cdot\omega + \varphi_V\cdot \eta$
    is an orientation form on $M$.
    By definition this form is nonzero
    on $U - U\cap V$ and $V - U\cap V$, and
    for all $p\in U\cap V$ and $e_1,\dots,e_n\in T_pM$,
    $\omega_p(e_1,\dots,e_n)$ has the same sign as $\eta_p(e_1,\dots,e_n)$,
    so $\varphi_U(p)\omega_p(e_1,\dots,e_n) +
    \varphi_V(p)\eta_p(e_1,\dots,e_n)$
    has the same sign as the values of $\omega$ and $\eta$,
    and in particular is nonzero.
\end{soln}

\begin{problem}
    [15-3]
    Suppose $n\ge 1$, and let $\alpha\colon S^n\to S^n$
    be the antipodal map $\alpha(x) = -x$.
    Show that $\alpha$ is orientation preserving
    if and only if $n$ is odd.
\end{problem}

\begin{soln}
    The vector $p\in S^n\subseteq\RR^{n+1}$
    is a unit normal to $T_pS^n$,
    so the normal vector field $X_p = p$
    is a nonvanishing section of the rank-$1$ normal bundle of $S^n$.
    The existence of this section tells us that
    $S^n$ has trivial normal bundle in $\RR^{n+1}$.

    Let $\omega$ be the volume form $\dd x_1\wedge\cdots\wedge\dd x_n$
    on $\RR^{n+1}$; by the solution to Problem 15-8(a) below,
    $\iota_X\omega$ is a volume form on $S^n$.
    Then $\alpha$ is orientation-preserving if and only if
    $\alpha^*(\iota_X\omega)$ induces the same orientation on $S^n$
    as $\iota_X\omega$
    (in general this holds for any diffeomorphism $\alpha$,
    since we want the differential at each point
    to induce an orientation preserving linear map
    on the top exterior product of each tangent space).

    Let $r_i\colon\RR^{n+1}\to\RR^{n+1}$ be the reflection map
    $r(x_1,\dots,x_i,\dots,x_n) = r(x_1,\dots,-x_i,\dots,x_n)$,
    so $\alpha$ is the composition of all the $r_i|_{S^n}$s;
    of course, each $r_i$ is a diffeomorphism
    and restricts to a diffeomorphism on $S^n$
    (we will identify $r_i$ and $r_i|_{S^n}$).
    Each $r_i$ reverses the orientation because
    for all $p\in S^n$ and $e_1,\dots,e_n\in T_pS^n$ a basis,
    \begin{equation*}
	(r_i^*(\iota_X\omega))_p(e_1,\dots,e_n)
	= \omega_p((\dd r_i)_pX_p, (\dd r_i)_pe_1,\dots,(\dd r_i)_pe_n)
	= -\omega_p(X_p,e_1,\dots,e_n)
	= -(\iota_X\omega)_p(e_1,\dots,e_n).
    \end{equation*}
    The last inequality is true because $\det(\dd r_i)_p = -1$
    and for a finite-dimensional vector space $V$
    and linear map $T\colon V\to V$,
    the induced map on the top exterior product space $\Lambda^{\dim V}V$
    is multiplication by the determinant.
    Finally, we find that
    \begin{equation*}
        \alpha^*(\iota_X\omega)
	= r_1^*\cdots r_{n+1}^*(\iota_X\omega)
	= (-1)^{n+1}\iota_X\omega,
    \end{equation*}
    so $\alpha$ is orientation-preserving
    if and only if $(-1)^{n+1}$ is positive,
    or equivalently $n$ is odd.
\end{soln}

\begin{problem}
    [15-5]
    Let $M$ be a smooth manifold with or without boundary.
    Show that the total spaces of $TM$ and $T^*M$ are orientable.
\end{problem}

\begin{soln}
    Let $p\in M$ be an arbitrary point, $U$ a neighborhood of $p$,
    and $(x_1,\dots,x_n)$ a local coordinate system
    in a neighborhood of $p$. These induce local coordinates
    $(x_1,\dots,x_n,v_1,\dots,v_n)$ on $\pi\inv(U)$.
    Define a local frame on $\pi\inv(U)$ by
    $(X^1_p, \dots, X^{2n}_p)
    = (\tanvec{x_1},\dots,\tanvec{x_n},\tanvec{x_1},\dots,\tanvec{x_n})$.
    If $\varphi\colon U\cap V\to U\cap V$ is the
    coordinate change map between two coordinate charts $U$ and $V$,
    then the corresponding coordinate change map
    between the charts on $U$ and $V$ has differential
    $\dd\varphi_p\oplus\dd\varphi_p$
    and Jacobian $\det(\dd\varphi_p)^2$,
    so the local frames on all charts defined in this way
    are consistenly oriented, implying $TM$ is orientable.

    Every smooth manifold admits a Riemannian metric,
    which then provides a natural isomorphism
    between the tangent and cotangent bundles
    (by the results of Chapter 13), so the cotangent bundle is orientable
    by virtue of being diffeomorphic to an orientable manifold.
\end{soln}

\begin{problem}
    [15-8]
    Suppose $M^m$ is an orientable Riemannian manifold
    and $S^s\subseteq M$ is an embedded or immersed submanifold.
    Prove that
    \begin{enumerate}[(a)]
        \ii If $S$ has trivial normal bundle, then $S$ is orientable.
	\ii If $S$ is an orientable hypersurface,
	then it has trivial normal bundle.
    \end{enumerate}
\end{problem}

\begin{soln}
    (a) Let $\omega$ be a volume form on $M$ and $i\colon S\to M$
    the immersion/embedding that defines $S$ as a submanifold of $M$.
    Since the normal bundle $N$ is trivial,
    it admits $r = \operatorname{rank} N$ linearly independent
    (global) sections $X_1$, \dots, $X_r$.
    Then, consider the $s$-form
    \[
	\eta = \omega(X_1,\dots,X_r, -)
	= \iota_{X_1}(\cdots \iota_{X_{r-1}}(\iota_{X_r}\omega))
    \]
    (where $\iota_V$ is interior multiplication by the vector field $V$).

    I claim that $i^*\eta$ is a volume form on $S$.
    To see this, we note that for all $p\in S$,
    if $e_1,\dots,e_s\in T_pS$ is a positively oriented basis of $T_pS$,
    then $di_pe_1$, \dots, $di_pe_s$ forms a basis
    of $T_{i(p)}M$ with the $(X_i)_p$
    because $i$ is an immersion, so
    \begin{equation*}
	(i^*\eta)_p(e_1,\dots,e_s)
	= \eta_{i(p)}(di_pe_1,\dots,di_pe_s)
	= \omega_{i(p)}((X_1)_p,\dots,(X_r)_p,
	di_pe_1,\dots,di_pe_s)\ne 0.
    \end{equation*}
    That is, $i^*\eta$ is nonzero
    at each point of $S$, so it defines an orientation on $S$.
    \\

    (b) For each $p\in S$, let $e_1,\dots,e_s\in T_pS$
    be a positively oriented basis of $S$
    (with respect to the orientation on $S$).
    Let $X_p$ be the unique unit normal vector at $p$
    such that $(X_p,e_1,\dots,e_s)$ is oriented positively
    according to the orientation on $M$.
    This is a smooth nonvanishing section of the normal bundle;
    smoothness follows from taking
    a positively oriented orthonormal frame around $p$
    which extends a positively oriented orthnormal frame for $S$ around $p$,
    since this necessarily forces the one vector
    not in $T_pS$ to equal $X_p$, and local frames
    consist of smooth vector fields.
    The existence of a nonvanishing global smooth section
    of the normal bundle (which has rank $1$) implies that it is trivial.
\end{soln}










\end{document}
